\begin{lemma}[Tashkinov augmentation sequences survive a component swap]
Let \(G\) be a finite loopless multigraph with its individual edge type,
let \(e_0\) be the distinguished edge, and let \(\varphi\) be a partial
\(k\)-edge-coloring.  Fix distinct colors \(\alpha,\beta\), an incidence
node \(z\), and put
\[
 \sigma=\noderef{TwoColorComponentSwap}(G,\varphi,\alpha,\beta,z).
\]
If \(Y\) is obtained from \(X\) by a
\(\varphi\)-\noderef{TashkinovAugmentationSequence} and both
\(\alpha\) and \(\beta\) belong to
\(\noderef{TreeMissingColors}(\sigma,X)\), then the same exact ordered
tree-sequences and the same individual new-edge identities form a
\(\sigma\)-Tashkinov augmentation sequence from \(X\) to \(Y\).
\end{lemma}
\begin{proof}
Write the given sequence, in the reflexive-transitive sense of
\noderef{TashkinovAugmentationSequence}, as
\[
                 X=X_0,X_1,\ldots,X_m=Y,
\]
where each nonempty step is a
\(\varphi\)-\noderef{TashkinovAugmentation}.  We induct on \(m\), while
retaining the following two assertions at stage \(i\): \(X\) is an exact
ordered \noderef{TreeSequenceInitialSegment} of \(X_i\), and the displayed
chain through \(X_i\) is a \(\sigma\)-augmentation sequence.  Both
assertions are immediate at \(i=0\).

Suppose they hold at \(i<m\), and unpack the next augmentation.  It gives
\[
 |X_{i+1}|=|X_i|+1,
 \qquad X_i\preccurlyeq X_{i+1},
\]
together with the exact new last-edge identity
\(g=X_{i+1}.\mathit{edge}(|X_i|)\) and a color \(c\) such that
\[
                 \varphi(g)=c,
        \qquad c\in\noderef{TreeMissingColors}(\varphi,X_i).  \tag{1}
\]
The old initial-segment assertion followed by
\(X_i\preccurlyeq X_{i+1}\) gives
\(X\preccurlyeq X_{i+1}\), entry by entry, so it remains only to verify
the color and missing-witness clauses of the same one-step augmentation
under \(\sigma\).

First suppose \(c\notin\{\alpha,\beta\}\).  The definition of
\noderef{TwoColorComponentSwap} changes only values \(\alpha\) and
\(\beta\), hence \(\sigma(g)=\varphi(g)=c\).  Moreover it fixes, throughout
the graph, every edge whose value is \(c\).  Expanding
\noderef{TreeMissingColors} in (1) gives a vertex
\(v\in\noderef{TreeSequenceVertexSet}(X_i)\) with
\(c\in\noderef{MissingColors}(G,\varphi,v)\).  Since all incident
\(c\)-colored edges are fixed, \(c\) is missing at \(v\) under \(\sigma\)
if and only if it is missing there under \(\varphi\).  Thus
\(c\in\noderef{TreeMissingColors}(\sigma,X_i)\), and all clauses of the
same exact augmentation hold under \(\sigma\).

Now suppose \(c\in\{\alpha,\beta\}\).  If the edge-node
\(\operatorname{inr}g\) is outside the selected two-color component, the
component-swap definition leaves its value equal to \(c\).  If it is
inside, the definition transposes its old value, so the new value is the
other member of \(\{\alpha,\beta\}\).  In either subcase there is a color
\(d\in\{\alpha,\beta\}\) for which
\[
                            \sigma(g)=d.                       \tag{2}
\]
By the two hypotheses of the lemma, \(d\) has a missing witness
\(w\in V(X)\) under \(\sigma\).  The induction invariant
\(X\preccurlyeq X_i\), expanded through
\noderef{TreeSequenceInitialSegment} and
\noderef{TreeSequenceVertexSet}, carries the very same vertex entry into
\(V(X_i)\).  Missingness is a property of the fixed graph coloring at
that vertex and is independent of which tree-sequence currently contains
it.  Hence
\[
                   d\in\noderef{TreeMissingColors}(\sigma,X_i). \tag{3}
\]
Equations (2)--(3), together with the unchanged length and exact
initial-segment data, verify every clause of
\noderef{TashkinovAugmentation} for the step
\(X_i\to X_{i+1}\) under \(\sigma\).

In both cases the transported step uses the original individual edge
identity \(g\), and appending it proves the two induction assertions at
\(i+1\).  At \(i=m\), the transported finite chain is exactly a
\(\sigma\)-\noderef{TashkinovAugmentationSequence} from \(X\) to \(Y\),
as required.
\end{proof}
